\clearpage\fancyhead[L]{\small KOMET.}\hypertarget{native-343-lower}{}\label{native-343-lower}
\begin{center}\textbf{KOMET.}\par\vspace{4pt}\hrule\vspace{4pt}\fontsize{10}{12}\selectfont Monatliche Correspondenz. Band XXIV, S. 406\,..\,411, Oktober 1811.\par\vspace{4pt}\hrule\vspace{4pt}\end{center}

\par\indent Die im vorigen Heft geäußerte Hoffnung, dass wir unseren Lesern in diesem Heft etwas Bestimmtes über die Umlaufzeit des Kometen mitteilen könnten, hat sich nicht erfüllt. Anfangs schien es wahrscheinlich, dass die starke, immer weiter zunehmende Abweichung der Beobachtungen von den parabolischen Elementen auf einer tatsächlichen Abweichung der Bahn von der Parabel beruhte. Später zeigte sich jedoch, dass eine kleine Korrektur der Elemente völlig ausreichte, um beides in befriedigende Übereinstimmung zu bringen. Da es von Interesse ist zu sehen, wie man die Elemente durch aufeinanderfolgende Änderungen immer besser mit dem Himmel in Übereinstimmung bringen kann, legen wir heute zunächst eine Tabelle vor, in der sämtliche bis zum 11. Oktober gemachten Beobachtungen mit den erstmals verbesserten Elementen von \textsc{Gauß} verglichen werden (Mon. Corresp., Septemberheft, S. 305) [S. 335 d. B.].

\Needspace{41\baselineskip}\SourceLeaf{344}
Vergleich sämtlicher Kometenbeobachtungen mit den parabolischen Elementen von \textsc{Gauß}.

\begin{center}

\fontsize{10}{12}\selectfont\setlength{\tabcolsep}{7pt}\renewcommand{\arraystretch}{1.05}\begin{tabular}{ll|c|c|l}
\multicolumn{2}{c|}{1811} & \multicolumn{2}{c|}{Abweichung} & Name \\
\multicolumn{2}{c|}{Beobachtungstag} & in Rektasz. & in Dekl. & des Beobachters \\\hline
August & 22 & $+61''$ & $-57''$ & Bessel \\
 & 23 & $+70$ & $-79$ & Bessel \\
 & 23 & $+119$ & $-21$ & Olbers \\
 & 25 & $+116$ & $-101$ & Olbers \\
 & 26 & $+98$ & $-84$ & Olbers \\
 & 27 & $+99$ & $-99$ & Olbers \\
 & 27 & $+63$ & $-50$ & Bessel \\
 & 28 & $+34$ & $-68$ & Bessel \\
 & 28 & $+46$ & $-124$ & Olbers \\
 & 29 & $+128$ & $-26$ & Oriani \\
 & 30 & $+9$ & $-41$ & Olbers \\
 & 30 & $+43$ & $-55$ & Bessel \\
 & 31 & $+37$ & $-49$ & Oriani \\
 & 31 & $+47$ & $-120$ & Olbers \\
Sept. & 1 & $+25$ & $-43$ & Oriani \\
 & 2 & $+19$ & $-74$ & Oriani \\
 & 3 & $+10$ & $-92$ & Oriani \\
 & 3 & $+23$ & $-67$ & Olbers \\
 & 4 & $+25$ & $-73$ & Oriani \\
 & 4 & $+51$ & $-129$ & Olbers \\
 & 4 & $+30$ & $-9$ & Gauss \\
 & 6 & $-43$ & $-89$ & Oriani \\
 & 6 & $-7$ & $-50$ & Lindenau \\
 & 6 & $+9$ & $-11$ & Gauss \\
 & 7 & $-26$ & $-74$ & Lindenau \\
 & 7 & $-17$ & $-80$ & Bessel \\
 & 7 & $-54$ & $+20$ & Gauss \\
 & 8 & $-54$ & $-45$ & Oriani \\
 & 8 & $-41$ & $-44$ & Lindenau \\
 & 9 & $-61$ & $-38$ & Oriani \\
 & 9 & $-62$ & $-7$ & Lindenau \\
 & 9 & $-65$ & $-64$ & Gauss \\
\end{tabular}
\end{center}
\Needspace{38\baselineskip}
\begin{center}

\fontsize{10}{12}\selectfont\setlength{\tabcolsep}{7pt}\renewcommand{\arraystretch}{1.05}\begin{tabular}{ll|c|c|l}
\multicolumn{2}{c|}{1811} & \multicolumn{2}{c|}{Abweichung} & Name \\
\multicolumn{2}{c|}{Beobachtungstag} & in Rektasz. & in Dekl. & des Beobachters \\\hline
Sept. & 10 & $-116''$ & $-27$ & Lindenau \\
 & 10 & $-54$ & $-71$ & Schubert \\
 & 11 & $-70$ & $-73$ & Oriani \\
 & 11 & $-135$ & $-28$ & Lindenau \\
 & 11 & $-82$ & $-106$ & Bessel. \\
 & 12 & $-120$ & $-43$ & Oriani \\
 & 12 & $-37$ & $-78$ & Bessel \\
 & 12 & $-86$ & $-53$ & Schubert \\
 & 13 & $-111$ & $-118$ & Oriani \\
 & 13 & $-130$ & $-34$ & Lindenau \\
 & 13 & $-119$ & $-74$ & Schubert \\
 & 14 & $-135$ & $-94$ & Oriani \\
 & 14 & $-146$ & $-34$ & Lindenau \\
 & 14 & $-140$ & .\;.\;. & Gauss \\
 & 14 & $-139$ & $-30$ & Schubert \\
 & 15 & $-154$ & $-51$ & Lindenau \\
 & 15 & $-195$ & $-76$ & Gauss \\
 & 15 & $-158$ & $-65$ & Schubert \\
 & 16 & $-225$ & $-54$ & Lindenau \\
 & 16 & $-219$ & $-30$ & Gauss \\
 & 16 & $-184$ & $-52$ & Schubert \\
 & 17 & $-172$ & $-43$ & Oriani \\
 & 17 & $-261$ & $-29$ & Lindenau \\
 & 18 & $-247$ & $-34$ & Lindenau \\
 & 19 & $-287$ & $-66$ & Lindenau \\
 & 22 & $-443$ & .\;.\;. & Bessel \\
 & 23 & $-420$ & .\;.\;. & Bessel \\
 & 26 & $-626$ & $+53$ & Lindenau \\
 & 30 & $-839$ & $+130$ & Lindenau \\
Okt. & 2 & $-879$ & $+118$ & Gauss \\
 & 5 & $-1123$ & $+336$ & Lindenau \\
 & 11 & $-1271$ & $+626$ & Lindenau \\
\end{tabular}

\end{center}

\Needspace{17\baselineskip}\par\indent Herr Professor \textsc{Gauß}, der uns in den letzten Tagen des Septembers und Anfang Oktober das große Vergnügen eines 14-tägigen Aufenthalts auf der hiesigen Sternwarte bereitete, war sehr geneigt, in diesen Abweichungen die Spur einer elliptischen oder hyperbolischen Bahn zu sehen. Er begann schon hier mit einigen Rechnungen zu ihrer Bestimmung, ohne jedoch zu einem bestimmten Ergebnis zu gelangen. Aber schon am zweiten Tag nach seiner Rückkehr nach Göttingen teilte er uns seine interessanten Ergebnisse mit. Wir geben sie hier mit seinen eigenen Worten wieder. „Datiert Göttingen, 14. Oktober. Ich habe meine letzten Rechnungen zum Kometen sofort überprüft und, wie ich zuvor vermutet hatte, darin bald einen Schreibfehler von 10 Minuten entdeckt. Nach dessen Korrektur fiel das Ergebnis ganz anders aus, und zwar zeigte sich,

\begin{center}
\emph{dass noch gar keine sichere Spur einer Elliptizität der Bahn zu bemerken ist.}
\end{center}

\noindent Meine verbesserten parabolischen Elemente lauten wie folgt.

\Needspace{12\baselineskip}
\begin{center}\fontsize{10}{12}\selectfont\setlength{\tabcolsep}{4pt}\renewcommand{\arraystretch}{1.15}\begin{tabular}{l}
Durchgang durch das Perihel 1811 Sept. 12\\
\quad $5^{\mathrm u}\,32^{\mathrm m}\,46^{\mathrm s}$ M. Z. in Göttingen\\
\end{tabular}\end{center}
\begin{center}\fontsize{10}{12}\selectfont\setlength{\tabcolsep}{3pt}\renewcommand{\arraystretch}{1.15}\begin{tabular}{llr}
Länge des Perihels & .\;.\;.\;.\;.\;.\;. & $75^\circ\,4'\,43''$\\
Log. des Abstands im Perihel & .\;.\;. & $0.015530$\\
Länge des aufsteigenden Knotens & . & $140^\circ\,21'\,40''$\\
Neigung der Bahn & .\;.\;.\;.\;.\;. & $73\;4\;18$\\
\end{tabular}\end{center}
\noindent Bewegung rückläufig

\Needspace{17\baselineskip}\SourceLeaf{345}
\par\indent Für die auf den Äquator bezogenen Koordinaten finde ich die folgenden Konstanten.

\[
\begin{array}{ll}
x = \dfrac{\alpha.\sin(v+348^\circ\,50'\,38'')}{\cos\frac{1}{2}v} & \log \alpha = 9.91435 \\[2ex]
y = \dfrac{\ObservedBeta.\sin(v+171^\circ\,56'\,51'')}{\cos\frac{1}{2}v} & \log \ObservedBeta = 9.80148 \\[2ex]
z = \dfrac{\gamma.\sin(v+80^\circ\;0'\;3'')}{\cos\frac{1}{2}v} & \log \gamma = 0.01538 \\
\end{array}
\]

\par\indent Diese Elemente sind übrigens nur flüchtig berechnet und ließen sich leicht noch besser an die Beobachtungen anpassen. Ich hielt es aber nicht für der Mühe wert, das schon jetzt zu tun. Sie reichen aus, um die Beobachtungen noch einige Zeit mit ihnen zu vergleichen, und zeigen, dass sich über die Elliptizität der Bahn bisher nichts Sicheres sagen lässt. Ich glaube nicht, dass die Umlaufzeit, wenn die Bahn elliptisch ist, unter 1000 Jahren liegen kann. Ich habe diese Elemente soeben mit Ihren und meinen Beobachtungen verglichen. Das bestätigt, was ich eben gesagt habe.

\Needspace{25\baselineskip}
\begin{center}\fontsize{10}{12}\selectfont\setlength{\tabcolsep}{6pt}\renewcommand{\arraystretch}{1.15}\begin{tabular}{l|c|c|c|c}
 & \multicolumn{4}{c}{Unterschied}\\
 & \multicolumn{2}{c|}{in Rektaszension} & \multicolumn{2}{c}{in Deklination} \\
1811 & v. L. & G. & v. L. & G. \\\hline
Sept. 4 & --- & $+5''$ & --- & $+23''$ \\
6 & --- & $-7$ & --- & $+28$ \\
7 & --- & $-58$ & --- & $+61$ \\
9 & --- & $-38$ & --- & $-18$ \\
10 & $-69''$ & & $+26''$ & \\
11 & $-70$ & & $+39$ & \\
14 & $-12$ & $-9$ & $+26$ & \\
15 & $+5$ & $-17$ & $+7$ & $-18$ \\
16 & $-37$ & $-36$ & $+10$ & $+35$ \\
17 & $-43$ & & $+33$ & \\
18 & $+4$ & & $+16$ & \\
19 & $-27$ & & $+1$ & \\
26 & $-47$ & & $+63$ & \\
30 & --- & $-50$ & $+76$ & \\
Oktober 2 & --- & $-29$ & & $+19$ \\
\end{tabular}
\end{center}

\par\indent Die Elemente ergeben also --- anscheinend während des gesamten Zeitraums dieser Beobachtungen --- um etwa eine halbe Minute zu kleine Rektaszensionen und ungefähr ebenso viel zu große Deklinationen. Dieser Unterschied ließe sich leicht beseitigen, ohne die Parabel zu verlassen und ohne die Übereinstimmung mit den Beobachtungen von \textsc{Zach} zu verschlechtern. \textsc{Olbers} hat mir in einem Brief, den ich hier vorfand, kurz von seiner einfallsreichen Hypothese über den Kometenschweif berichtet. Er nimmt einen Stoff an, der vom Kometen erzeugt und sowohl vom Kometen als auch von der Sonne abgestoßen wird. Der Stoff häuft sich dort an, wo die beiden Abstoßungskräfte eine Art Gleichgewicht halten, und bildet eine Art hohlen parabolischen Konoid. \textsc{Olbers} wird Ihnen darüber ausführlicher schreiben.
